\section{Introduction}
% TODO: port references.bib dari docs/PAPER.md (cite keys di bawah mengacu ke sana).

Generative autoregressive language models are increasingly deployed for
structured enterprise decisions: payment fraud routing, customer inquiry
triaging, and policy enforcement. In this setup the model is prompted to emit
its decision as a formatted payload such as JSON or XML. This \textit{System
Two} paradigm carries three operational costs. First, latency: emitting a
formatted decision requires $O(T)$ sequential decoding steps, spending GPU
cycles on boilerplate punctuation while holding the key-value cache. Second,
syntactic reliability: even with constrained decoding, autoregressive sampling
retains a non-zero probability of malformed output. Third, epistemic
pathology: preference-based post-training rewards confident answers, so raw
output scores are systematically overconfident and unusable for risk-controlled
automation~\cite{guo2017calibration}.

\textit{System One} decision models separate structured judgment from text
generation: an input state $S$ and a typed question schema $Q$ are mapped
directly to a joint probability distribution over candidate options in a
single forward pass ($O(1)$ decoding steps), with no free-form text and no
output parsing. The paradigm was pioneered by TypeSafe Jev (closed API),
Cloudflare Clef~\cite{cloudflare2025clef}, Convai Laya, and the open Unsloth
training recipe~\cite{unsloth2026decision}. Existing work concentrates on dense
bidirectional encoders or decoder-only transformers; hybrid
convolution-attention backbones at sub-billion scale, attractive for edge
deployment, remain an unexplored substrate for decision heads. We study
LiquidAI LFM2.5-350M~\cite{liquidai2025lfm}, a 350M hybrid model combining
gated-convolution blocks with grouped-query attention.

Rather than proposing a new architecture (our decision head follows the Clef
design), this paper asks which \textit{training recipe} choices make small
decision models transfer: (RQ1) does schema augmentation (option renaming,
description paraphrasing, derived questions) produce transfer to unseen
decision taxonomies; (RQ2) how does data-mixture composition trade off across
task groups under fixed 350M capacity; with calibrated probabilities and
selective-classification coverage evaluated throughout.

Our contributions are: (1) an adaptation recipe and open artifact,
\texttt{itsdevice/lfm-2.5-350m-decs1}; (2) controlled experiments ---
multi-seed runs, leave-one-taxonomy-out evaluation, and token-matched
augmentation ablations [PENDING run \#1--8, tracker
docs/RECOMMENDATION.md \S6.4]; (3) a calibration evaluation (proper scoring
rules plus post-hoc temperature scaling) with selective coverage at fixed
risk; and (4) a capacity-frontier analysis against a 0.8B dense baseline run
under the identical protocol [PENDING run \#8]. This is an empirical systems
paper: every comparative claim traces to a logged run, and result numbers
enter this section only after the tracker turns green.

The remainder of the paper covers related work, the architecture and training
methodology, the dataset composition and augmentation pipeline, the
experimental results, and deployment discussion with limitations.
